problem:  
Let \((M,\mathcal{D},g,\mu)\) be a compact, equiregular sub‑Riemannian manifold with Hausdorff dimension \(Q>n=\dim M\).  
For each \(\varepsilon>0\), consider the Riemannian approximation  
\[
g_\varepsilon(v)=g(v_\mathcal{D})+\varepsilon^{-2}g(v_{\mathcal{D}^\perp}),
\]
and the associated metric measure spaces \((M,d_{g_\varepsilon},\mu)\), so that \(d_{g_\varepsilon}\to d_{SR}\) in measured Gromov–Hausdorff sense as \(\varepsilon\to0\).  

Suppose—hypothetically—that there exists a sequence \(\varepsilon_j\to0\) such that each \((M,d_{g_{\varepsilon_j}},\mu)\) satisfies the uniform synthetic curvature–dimension bound \(\mathrm{RCD}^*(K_0,N_0)\).  

**Question:**  
Is it possible to formulate a synthetic curvature–dimension structure, denote it \(\mathrm{RCD}^*_{\mathrm{subR}}(K_0,N_0)\), on the limit space \((M,d_{SR},\mu)\) satisfying the following explicit properties:

1. (**Consistency**) If tangent cones are Euclidean, it reduces to the standard \(\mathrm{RCD}^*(K_0,N_0)\) structure.  
2. (**Non‑Hilbertian stability**) The condition is stable under measured Gromov–Hausdorff convergence of such degenerating families \( (M,d_{g_\varepsilon},\mu)\to (M,d_{SR},\mu) \).  
3. (**Analytic formulation**) There exists a well‑defined **graded Cheeger energy** \({\rm Ch}^{\rm grad}\) on \( (M,d_{SR},\mu)\), quadratic on the horizontal component but possibly non‑quadratic on the vertical, satisfying an “entropy convexity” property along subunit Wasserstein geodesics.  
4. (**Tangent structure**) For a.e. \(x\), the measured tangent cone coincides with a Carnot group equipped with its Popp measure; the new curvature notion should encode this algebraic geometry.  

Determine whether these four axioms can coexist without contradiction.  
If they can, construct such a \( \mathrm{RCD}^*_{\mathrm{subR}}(K,N) \) framework; if they cannot, prove that at least one of (1)–(4) is impossible (for instance, show that entropy convexity forces Hilbertianity of the tangent module, thereby precluding Carnot tangents).  

Why is it a "good" problem:  
This refined formulation transforms a conceptual program into a **precisely testable geometric dilemma**: can the desirable analytic and geometric properties of the \(\mathrm{RCD}^*\) theory survive the loss of Hilbertian infinitesimal structure inherent in sub‑Riemannian limits?  

It explicitly ties together geometry (Carnot tangent cones), analysis (Cheeger energies and modules), and optimal transport (entropy convexity), revealing the fundamental incompatibility or potential synthesis among them. The outcome—either a successful construction or a rigidity/impossibility result—would settle a central question about the reach of synthetic curvature theory.  

The problem remains simple in statement yet vast in implication: one must either design or obstruct a coherent “sub‑Hilbertian curvature theory” that unifies Riemannian and sub‑Riemannian geometries. Its resolution would profoundly extend our understanding of curvature, measure, and transport, thus perfectly fulfilling the criteria of a “good” mathematical problem.